\documentclass[10pt,a4paper]{ctexart}
\usepackage[margin=1.35cm]{geometry}
\usepackage{amsmath,amssymb,array,booktabs,multicol}
\setlength{\parindent}{0pt}
\setlength{\parskip}{2pt}
\pagestyle{plain}

\begin{document}
\begin{center}
  {\LARGE 第五章 ODE 上机公式速查}\\
  \small 只保留可直接写进答案和代码的公式
\end{center}

\begin{multicols}{2}
\section*{0. 通用写法}
\[
y'=f(t,y),\quad y(a)=y_0,\qquad t_n=a+nh.
\]
单步法统一写为
\[
y_{n+1}=y_n+h\Phi(t_n,y_n,h).
\]
若已知精确解，计算
\[
E(h)=\max_n\|y(t_n)-y_n\|_\infty,\qquad
p\approx\frac{\log(E(h)/E(h/2))}{\log 2}.
\]

\section*{1. Euler 与二阶 RK}
\textbf{向前 Euler（一阶）}
\[
y_{n+1}=y_n+h f(t_n,y_n).
\]

\textbf{改进 Euler / 显式梯形（二阶）}
\[
K_1=f(t_n,y_n),\quad
K_2=f(t_n+h,y_n+hK_1),
\]
\[
y_{n+1}=y_n+\frac h2(K_1+K_2).
\]

\textbf{中点 RK2 / 变形 Euler（二阶）}
\[
K_1=f(t_n,y_n),\quad
K_2=f\left(t_n+\frac h2,y_n+\frac h2K_1\right),
\quad y_{n+1}=y_n+hK_2.
\]

\textbf{教材二阶 Heun（二阶）}
\[
K_1=f(t_n,y_n),\quad
K_2=f\left(t_n+\frac{2h}{3},y_n+\frac{2h}{3}K_1\right),
\]
\[
y_{n+1}=y_n+\frac h4(K_1+3K_2).
\]

\section*{2. 经典 RK4}
\[
\begin{aligned}
K_1&=f(t_n,y_n),\\
K_2&=f(t_n+h/2,y_n+hK_1/2),\\
K_3&=f(t_n+h/2,y_n+hK_2/2),\\
K_4&=f(t_n+h,y_n+hK_3),\\
y_{n+1}&=y_n+\frac h6(K_1+2K_2+2K_3+K_4).
\end{aligned}
\]
\textbf{RK4 Butcher 表}
\[
\begin{array}{c|cccc}
0\\
\frac12&\frac12\\
\frac12&0&\frac12\\
1&0&0&1\\ \hline
&\frac16&\frac13&\frac13&\frac16
\end{array}
\]

\section*{3. 一般显式 RK 与 Butcher 表}
\[
K_i=f\left(t_n+c_i h,\ y_n+h\sum_{j<i}a_{ij}K_j\right),
\quad
y_{n+1}=y_n+h\sum_{i=1}^s b_iK_i.
\]
\[
\begin{array}{c|c}
\boldsymbol c&A\\ \hline
&\boldsymbol b^T
\end{array}
\]

\section*{4. RKF45 自适应}
\textbf{六个斜率}
\[
\begin{aligned}
K_1={}&f(t,y),\\
K_2={}&f(t+\tfrac14h,\ y+h\tfrac14K_1),\\
K_3={}&f(t+\tfrac38h,\ y+h(\tfrac3{32}K_1+\tfrac9{32}K_2)),\\
K_4={}&f(t+\tfrac{12}{13}h,\ y+h(\tfrac{1932}{2197}K_1-\tfrac{7200}{2197}K_2+\tfrac{7296}{2197}K_3)),\\
K_5={}&f(t+h,\ y+h(\tfrac{439}{216}K_1-8K_2+\tfrac{3680}{513}K_3-\tfrac{845}{4104}K_4)),\\
K_6={}&f(t+\tfrac12h,\ y+h(-\tfrac8{27}K_1+2K_2-\tfrac{3544}{2565}K_3+\tfrac{1859}{4104}K_4-\tfrac{11}{40}K_5)).
\end{aligned}
\]

\textbf{四阶、五阶解与误差}
\[
y_4=y+h\left(\frac{25}{216}K_1+\frac{1408}{2565}K_3+
\frac{2197}{4104}K_4-\frac15K_5\right),
\]
\[
y_5=y+h\left(\frac{16}{135}K_1+\frac{6656}{12825}K_3+
\frac{28561}{56430}K_4-\frac9{50}K_5+\frac2{55}K_6\right),
\]
\[
\mathrm{err}=\|y_5-y_4\|_\infty.
\]
\[
\mathrm{err}\leq\mathrm{tol}:\ \text{接受， }t\leftarrow t+h,\ y\leftarrow y_5;
\]
\[
\mathrm{err}>\mathrm{tol}:\ \text{拒绝，不更新 }t,y.
\]
\[
h_{\rm new}=0.84h\left(\frac{\mathrm{tol}}{\mathrm{err}}\right)^{1/4},
\quad h_{\min}\leq h_{\rm new}\leq h_{\max}.
\]
若 $\mathrm{err}=0$，直接放大步长；每步先取
\[
h\leftarrow\min(h,T-t)
\]
保证最后一步不越过终点。

\textbf{RKF45 Butcher 表}
\[
\renewcommand{\arraystretch}{1.15}
\begin{array}{c|rrrrrr}
0\\
\frac14&\frac14\\
\frac38&\frac3{32}&\frac9{32}\\
\frac{12}{13}&\frac{1932}{2197}&-\frac{7200}{2197}&\frac{7296}{2197}\\
1&\frac{439}{216}&-8&\frac{3680}{513}&-\frac{845}{4104}\\
\frac12&-\frac8{27}&2&-\frac{3544}{2565}&\frac{1859}{4104}&-\frac{11}{40}\\ \hline
y_4&\frac{25}{216}&0&\frac{1408}{2565}&\frac{2197}{4104}&-\frac15&0\\
y_5&\frac{16}{135}&0&\frac{6656}{12825}&\frac{28561}{56430}&-\frac9{50}&\frac2{55}
\end{array}
\]

\section*{5. 方程组与高阶方程}
向量方程直接套同一个 RK 公式：
\[
\boldsymbol y'=\boldsymbol f(t,\boldsymbol y),\quad
\boldsymbol y(a)=\boldsymbol y_0.
\]
二阶方程
\[
y''=g(t,y,y'),\quad y(a)=\eta_1,\ y'(a)=\eta_2
\]
令 $w_1=y,\ w_2=y'$，则
\[
\begin{pmatrix}w_1'\\w_2'\end{pmatrix}
=
\begin{pmatrix}w_2\\g(t,w_1,w_2)\end{pmatrix},
\quad
\boldsymbol w(a)=\begin{pmatrix}\eta_1\\\eta_2\end{pmatrix}.
\]
$m$ 阶方程：令 $w_1=y,w_2=y',\ldots,w_m=y^{(m-1)}$。

\section*{6. Adams 最低限度模板}
记 $f_n=f(t_n,y_n)$。四步 AB 预报：
\[
y_{n+1}^{P}=y_n+\frac h{24}
(55f_n-59f_{n-1}+37f_{n-2}-9f_{n-3}).
\]
三步 AM 校正：
\[
y_{n+1}=y_n+\frac h{24}
(9f_{n+1}^{P}+19f_n-5f_{n-1}+f_{n-2}).
\]
PECE：Predict $\to$ Evaluate $\to$ Correct $\to$ Evaluate。\\
必须先用 RK4 得到 $y_1,y_2,y_3$；Adams 必须固定步长。

\section*{7. 误差与阶数}
若方法为 $p$ 阶：
\[
\text{局部截断误差}=O(h^{p+1}),\qquad
\text{整体误差}=O(h^p).
\]
\[
p\approx
\frac{\log(E(h_1)/E(h_2))}
{\log(h_1/h_2)}.
\]
常见整体阶数：
\[
\begin{array}{c|ccccc}
\text{方法}&\text{Euler}&\text{改进Euler}&\text{中点}&\text{Heun}&\text{RK4}\\ \hline
p&1&2&2&2&4
\end{array}
\]

\section*{8. 刚性题最低限度}
测试方程 $y'=\lambda y$：
\[
\text{Euler 稳定： }-2<h\lambda<0.
\]
线性系统 $\boldsymbol y'=A\boldsymbol y+\boldsymbol g(t)$ 的隐式梯形：
\[
\left(I-\frac h2A\right)\boldsymbol y_{n+1}
=\left(I+\frac h2A\right)\boldsymbol y_n
+\frac h2[\boldsymbol g(t_n)+\boldsymbol g(t_{n+1})].
\]
若显式方法突然发散，先检查步长是否超出稳定范围。

\section*{9. 考试固定答案句式}
\begin{enumerate}
\item 设步长 $h=(T-a)/N$，令 $t_n=a+nh$，取 $y_0=y(a)$。
\item 按所选方法计算各斜率并递推得到 $y_{n+1}$。
\item 若为高阶方程，先令 $w_1=y,w_2=y',\ldots$ 化为一阶系统。
\item 若有精确解，列表给出 $t_n,y_n,y(t_n),|y(t_n)-y_n|$，并计算最大误差。
\item 用 $h,h/2,h/4$ 重算，按对数误差比估计收敛阶。
\item 自适应法中，误差合格才更新时间和解；最后一步缩短至 $T$。
\end{enumerate}

\end{multicols}
\end{document}
